% !TEX root = IE3_expE_report.tex
\documentclass[lualatex,a4paper,10pt]{jlreq}
% Shared LuaLaTeX preamble.
\usepackage{luatexja-fontspec}
\setmainfont{TeX Gyre Termes}
\setsansfont{TeX Gyre Heros}
\setmonofont{Latin Modern Mono}
\setmainjfont[BoldFont=HaranoAjiGothic-Medium]{HaranoAjiMincho-Regular}
\setsansjfont[BoldFont=HaranoAjiGothic-Bold]{HaranoAjiGothic-Medium}
\usepackage[top=22mm,bottom=22mm,left=23mm,right=23mm]{geometry}
\usepackage{amsmath,amssymb,booktabs,array,longtable,tabularx}
\usepackage{graphicx,xcolor,tikz,circuitikz}
\usetikzlibrary{arrows.meta,positioning,calc}
\usepackage{enumitem,fancyhdr,listings,needspace}
\usepackage[unicode,colorlinks=true,linkcolor=labblue,urlcolor=labteal]{hyperref}
\usepackage{bookmark}
\definecolor{labblue}{RGB}{30,70,112}
\definecolor{labteal}{RGB}{0,100,105}
\definecolor{labgray}{gray}{0.95}
\ModifyHeading{section}{font={\large\sffamily\bfseries\color{labblue}}}
\ModifyHeading{subsection}{font={\normalsize\sffamily\bfseries\color{labteal}}}
\setlength{\parindent}{1\zw}
\setlength{\parskip}{3pt}
\setlength{\emergencystretch}{2em}
\renewcommand{\arraystretch}{1.35}
\setlist{itemsep=3pt,topsep=4pt,leftmargin=2\zw}
\newcolumntype{Y}{>{\raggedright\arraybackslash}X}
\newcolumntype{C}[1]{>{\centering\arraybackslash}p{#1}}
\pagestyle{fancy}
\fancyhf{}
\fancyhead[L]{\small\color{labblue} IE3 後期・電子工学実験}
\fancyhead[R]{\small テーマ\LabCode}
\fancyfoot[C]{\thepage}
\setlength{\headheight}{16pt}
\DeclareOldFontCommand{\rm}{\normalfont\rmfamily}{\mathrm}
\lstset{basicstyle=\small\ttfamily,columns=fullflexible,keepspaces=true,
 breaklines=true,frame=single,backgroundcolor=\color{labgray},
 numbers=left,numberstyle=\tiny,showstringspaces=false,aboveskip=8pt,belowskip=8pt}
\newcommand{\blank}{\rule{22mm}{0.3pt}}
\newcommand{\answerline}{\par\noindent\rule{\linewidth}{0.25pt}\par}
\newcommand{\writing}[1]{\par\Needspace{\dimexpr#1+5mm\relax}\noindent%
 \fcolorbox{labblue!35}{white}{\parbox[c][#1][t]{\dimexpr\linewidth-2\fboxsep-2\fboxrule\relax}{\vspace{2mm}\noindent\strut\hfill\par}}\par}
\newcommand{\hint}[1]{\par\smallskip\noindent{\sffamily\bfseries\color{labteal}記述の視点：}#1\par}
\newcommand{\note}[1]{\par\smallskip\noindent\fcolorbox{labblue!45}{labblue!5}{%
 \parbox{\dimexpr\linewidth-2\fboxsep-2\fboxrule\relax}{#1}}\par\smallskip}
\newcommand{\eqerror}{\varepsilon=\frac{x_{\mathrm{meas}}-x_{\mathrm{th}}}{x_{\mathrm{th}}}\times100\ \%}
\newcommand{\LabSource}{徳山工業高等専門学校 情報電子工学科，
 『2026年度 電子工学実験（3年後期）テキスト』，テーマ\LabCode，
 \expandafter\path\expandafter{\SourcePdf}。}
\newcommand{\reporttitle}{
 \noindent{\small\sffamily\color{labteal}実験報告書・記入ガイド}\par
 \noindent{\LARGE\sffamily\bfseries \LabCode\quad\LabTitle}\par\smallskip
 \noindent 氏名：\blank\quad 班：\blank\quad 実験日：\blank\par
 \note{目的・原理・手順を確認し、実測値と自分の考察を記入する。
 考察のガイドは、検討する事象・使う測定結果・記述の方針を示す。}
}
\newcommand{\prelabtitle}{
 \noindent{\small\sffamily\color{labteal}事前学習・前提知識プリント}\par
 \noindent{\LARGE\sffamily\bfseries \LabCode\quad\LabTitle}\par\smallskip
 \noindent 氏名：\blank\quad 班：\blank\quad 予習日：\blank\par
 \note{用語、回路の動き、計算、測定表への記入の順に読む。
 例題の値は説明用であり、実験結果ではない。}
}
\newcommand{\tocrow}[3]{\hyperref[#1]{#2}& #3 &\pageref*{#1}\\[2mm]}
\newcommand{\documentmap}[1]{
 \section*{目次と概要}
 \noindent 青色の項目をクリックすると該当箇所へ移動する。\par
 {\small\begin{longtable}{@{}p{0.29\linewidth}p{0.61\linewidth}r@{}}
 \toprule {\color{labblue}\bfseries 項目} & {\color{labblue}\bfseries 読むと分かること} & 頁\\\midrule
 \endhead #1\bottomrule
 \end{longtable}}
 \clearpage
}
\newenvironment{terms}{\par\smallskip\begin{longtable}{@{}p{0.23\linewidth}p{0.73\linewidth}@{}}
 \toprule {\color{labteal}\bfseries 用語}&{\color{labteal}\bfseries 意味・回路での役割}\\\midrule\endhead}
 {\bottomrule\end{longtable}}
\newcommand{\term}[2]{\textbf{#1}&#2\\[2mm]}
\newcommand{\guidepoint}[4]{
 \Needspace{32mm}\subsection{#1}
 \noindent{\bfseries\color{labteal}考える事象：}#2\par
 \noindent{\bfseries\color{labteal}参照する結果：}#3\par
 \noindent{\bfseries\color{labteal}記述の方針：}#4\par
}
\newcommand{\reportclosing}[1]{
 \clearpage\section{結論のまとめ方}\label{sec:closing}
 \hint{#1}
 \ConclusionGuide
 \begin{enumerate}
 \item 目的に対応する最も重要な実測結果を、測定条件と数値を添えて述べる。
 \item 原理と一致した点・一致しなかった点を示す。原因は確認済みの事実と推測を分ける。
 \item 実施範囲で言えることと、未確認の条件を一文ずつまとめる。
 \end{enumerate}
 \note{書き出しの型：「○○の条件で△△を測定したところ、□□となった。
 これは◇◇と整合する。ただし××は確認しておらず、一般化には追加測定が必要である。」
 空欄は自分の実測値と根拠で埋める。}
 \noindent 結論（実験後に記入）：\writing{30mm}
 \noindent 改善案・次に確認したい条件：\writing{16mm}
 \section*{参考文献}\label{sec:references}
 \begin{enumerate}
 \item \LabSource
 \item 使用したデータシート・書籍・ウェブ資料（著者、題名、該当箇所、URL、参照日）を追加する。
 \end{enumerate}
}
\newcommand{\prelabclosing}{
 \section*{準備・出典}\label{sec:prelabsource}
 \begin{itemize}[nosep]
 \item 資料、筆記具、記録用紙、電卓と、実験条件・機器設定の記録欄。
 \item 確認問題の解答と、分からない用語・回路点への具体的な質問。
 \end{itemize}
 {\small\noindent\LabSource\par}
}
\newcommand{\simpleflow}[1]{
 \begin{center}
 \begin{tikzpicture}[node distance=5mm and 5mm,
 every node/.style={font=\small},box/.style={draw=labblue,fill=labblue!4,rounded corners=1mm,
 align=center,minimum height=10mm,text width=30mm}]
 #1
 \end{tikzpicture}
 \end{center}
}

\newcommand{\LabCode}{E}
\newcommand{\LabTitle}{マルチバイブレータ}
\newcommand{\SourcePdf}{IE3_expE_A4.pdf}
\newcommand{\ConclusionGuide}{離散回路と555それぞれについて、非安定は二状態の繰返し、単安定はトリガ後の復帰を実測波形で要約する。
主要な幅・周期の一致度と抵抗変更の傾向を示し、入力条件と使用した素子値を添える。\par}
\begin{document}
\reporttitle
\documentmap{
\tocrow{sec:auto1}{1 目的}{RC充放電で切り替わる非安定・単安定回路を実験で確かめる。}
\tocrow{sec:auto2}{2 原理}{RC充放電で切り替わる非安定・単安定回路の仕組みと成立条件を整理する。}
\tocrow{sec:auto3}{3 装置・条件}{機器・定数・測定点を確認し、資料の順序で操作する。}
\tocrow{sec:auto4}{4 方法}{機器・定数・測定点を確認し、資料の順序で操作する。}
\tocrow{sec:auto5}{5 結果}{節点の波形、パルス幅、周期を表や図に記録する。}
\tocrow{sec:auto6}{6 考察：結果を根拠に説明する}{節点の波形、パルス幅、周期を根拠に、比較と原因の検討を進める。}
\tocrow{sec:closing}{7 結論のまとめ方}{主要な結果、成立範囲、未確認の条件を短くまとめる。}
\tocrow{sec:references}{参考文献}{使用した資料と外部の根拠を示す。}
}

\section{目的}\label{sec:auto1}
トランジスタおよび555タイマによる非安定・単安定マルチバイブレータを動作させ、コンデンサの充放電と状態遷移の対応を確かめる。
\section{原理}\label{sec:auto2}
非安定回路は二状態を繰り返し、単安定回路はトリガ後に一定時間だけ状態を変えて元に戻る。
離散トランジスタ回路の近似時間は
\[
 T_1\simeq0.69C_1R_{B1},\quad T_2\simeq0.69C_2R_{B2},
 \quad T=T_1+T_2
\]
である。555では内部の$V_{CC}/3$と$2V_{CC}/3$のしきい値を用い、
\[
 t_H=\ln2(R_A+R_B)C,\quad t_L=\ln2R_BC,\quad
 f=\frac1{t_H+t_L},\quad t_{\rm mono}=\ln3R_AC
\]
となる。
\section{装置・条件}\label{sec:auto3}
トランジスタ非安定：$V_{CC}=10\,\mathrm{V}$。
単安定：$V_{CC}=+10\,\mathrm{V},V_{BB}=-10\,\mathrm{V}$。
555：$V_{CC}=5\,\mathrm{V}$。実際の素子定数と機器設定を記録する。
\writing{16mm}
\section{方法}\label{sec:auto4}
\begin{enumerate}
\item 非安定トランジスタ回路を組み、指定4点の波形を共通の基準波形と対比する。周期と各状態の持続時間を測る。
\item 単安定回路を組み、無入力時の安定状態を確認する。矩形波$10\,\mathrm{V_{pp}}$、オフセット$5\,\mathrm{V}$を与え、入力・微分波形を含む指定6点を観測する。
\item 555非安定回路で$R_A=5\,\mathrm{k}\Omega,R_B=3\,\mathrm{k}\Omega,C=0.1\,\mu\mathrm{F}$とし、高・低時間、周期を測る。
\item 555単安定回路で$R_A=5.6\,\mathrm{k}\Omega,C=0.1\,\mu\mathrm{F}$とし、資料指定の$1\,\mathrm{kHz}$、デューティ80\%のトリガを与える。実際の高・低電圧を記録する。
\item 電源を切って$R_A$を$10\,\mathrm{k}\Omega$可変抵抗に変更し、パルス幅の変化を調べる。
\end{enumerate}
\clearpage
\section{結果}\label{sec:auto5}
表1：時間特性。理論欄は使用した実際の素子値で計算する。
\par\noindent\begin{tabularx}{\linewidth}{YYYYY}
\toprule 回路 & 項目 & 理論/ms & 実測/ms & 誤差/\%\\\midrule
離散非安定 &$T_1$&&&\\ &$T_2$&&&\\ &$T$&&&\\
離散単安定 &幅&&&\\555非安定 &$t_H$&&&\\ &$t_L$&&&\\ &$T$&&&\\
555単安定 &幅&&&\\\bottomrule
\end{tabularx}
図1：離散非安定の4点、単安定の6点の波形（同一時間軸の図を別紙に追加する）。
\writing{40mm}
図2：555非安定・単安定の出力とコンデンサ電圧。しきい値を示す。
\writing{35mm}
\clearpage
表2：可変抵抗と単安定パルス幅。
\par\noindent\begin{tabularx}{\linewidth}{YYY}
\toprule 実際の$R_A$/$\mathrm{k}\Omega$ & 理論幅/ms & 実測幅/ms\\\midrule
&&\\&&\\&&\\&&\\\bottomrule
\end{tabularx}
\clearpage\section{考察：結果を根拠に説明する}\label{sec:auto6}
\guidepoint{離散回路の切替えと各節点波形}
{コレクタの急変とベース・容量の変化がどの順で起こるか。}
{非安定4点、単安定6点の波形、トリガと微分点、実際の回路図。}
{同じ時間軸へ切替え時刻を記し、各トランジスタのON/OFFを説明する。スピードアップ容量・保護ダイオードは接続先と観測点を対応させて役割を述べる。}
\guidepoint{555の幅・周期と理論値}
{充放電経路ごとの実測時間がRC式に近いか。}
{表1の$t_H,t_L,T$と単安定幅、実際のRとC、コンデンサ電圧。}
{非安定と単安定を分けて比較する。しきい値に達する時刻を波形で示し、容量公差・電源・負荷を差の候補とする。}
\guidepoint{抵抗変更で幅が変わるか}
{$R_A$変更に対して単安定幅が比例するか。}
{表2の抵抗実値と幅、同じ容量・トリガ条件。}
{幅対Rのグラフを作り、理論の傾き$\ln3C$と比較する。抵抗のつまみ位置ではなく実値を用い、入力条件をそろえる。}
\guidepoint{トリガの影響と測定の範囲}
{入力の幅・極性や観測負荷が動作へ影響するか。}
{入力の高低電圧・デューティ、出力幅、微分波形、プローブ設定。}
{入力と出力のデューティを混同しない。通常の幅公式から外れる条件では、長いトリガや再トリガの扱いを回路条件から検討する。未試験の条件は今後の確認として記す。}
\subsection{資料の課題・追加検討}
\begin{enumerate}
\item 各波形の立上り・立下りと、トランジスタのON/OFF、コンデンサ電圧の変化を対応させる。
\item 微分回路とスピードアップ用コンデンサの役割を説明する。保護ダイオードの働きも回路から検討する。
\item $R_{B1}=R_{B2}=75\,\mathrm{k}\Omega,T_1=1\,\mathrm{ms},T_2=0.5\,\mathrm{ms}$を実現する容量を計算する。
\item 比較器、ラッチ、放電トランジスタを使って555の動作を説明し、$\ln2,\ln3$の由来を充放電式から示す。
\item 抵抗・容量誤差、トリガ幅、測定点の負荷が時間に及ぼす影響を述べる。
\end{enumerate}
\reportclosing{非安定と単安定の違い、RC時定数と出力時間の関係を要約する。}
\end{document}
